\section{Introduction}

An LLM agent negotiating with a competitor must reason about more than the object being exchanged. A sale can finance the seller, complete the buyer's productive assets, and alter the position of a third party. A convincing explanation of those consequences is one form of evidence; a useful agreement and a correctly executed investment are others. Understanding economic agency requires connecting them while keeping their success criteria distinct.

Monopoly provides a compact persistent economy in which these links can be observed. Negotiated deeds become opportunities to build, development creates rent obligations, and bankruptcy transfers assets to a creditor. The player who pays for a completing deed may not be the player who ultimately earns its rent. Other participants can anticipate those consequences when pricing an asset, refusing a swap or deciding whether to keep a rival solvent. The game therefore connects language-mediated bargaining with checkable resource constraints and subsequent multi-player outcomes.

We study twenty complete games played by four configured frontier-model systems: GPT-6.1 Sol, Gemini 4 Argon, Claude Opus 5.5 and Grok 4.7. The records contain 6,989 decisions, 526 negotiation threads and every applied action. The study is observational. Its central question is how strategic attention expressed in bargaining relates to executable economic behavior, rather than which underlying model is best at Monopoly.

Three findings organize the paper. First, bargaining responds to what an exchange enables for another player. Refusals, counteroffers and recorded reasons reveal attention to rival development and liquidity. Asking-price differences persist within seller and mortgage categories, but the small set of same-deed comparisons is mixed. This identifies a behavioral pattern worth explaining without turning it into an estimated strategic premium. Second, agents execute many multi-step financing tasks: 128 of 188 episodes requiring liquidation end in payment, sixty-one through multiple actions. This positive census complements the smaller set of payable declarations and construction requests that do not execute. Third, the selected conduct trajectories connect bargaining to future ownership. Agents bluff about limits, selectively describe collateral, ask others to deny financing, and withdraw an assurance while continuing to bid for the same group.

Our contributions are an empirical account of relational bargaining using complete offer and counteroffer populations; a census of cash-shortfall resolution linked to development and estate flows; and eight source-grounded communication trajectories that distinguish what was said, received and done. The artifact trail supplies the method supporting these findings. It is not evidence that a new general auditor has been validated. Essential prompt and response conditions are reported before the results, including provider refusals that change one participant's private-note and retry opportunities.

The motivating questions span strategic reasoning, opponent modeling, negotiation and deception, planning, adaptation, economic risk, collective behavior and evaluation validity. We preserve this breadth as a coverage map, not as a claim that twenty games measure every construct. The most direct evidence concerns stated strategic considerations and local execution. General-versus-multi-agent capability, causal learning, persuasion and unfamiliar-opponent transfer require a different design. This distinction lets the paper explain what the agents do without asking outcome counts to answer questions they cannot identify.
